% !TEX root = ../b-rg.tex
\chapter{Orthography}\label{ch:orthography}

\section{The alphabet}\label{sec:alphabet}
The primary letters in the Borlish alphabet are ordered as in the table below.
\begin{table}[ht]
    \centering
    \begin{tabular}{c c | c c | c c}
        \multicolumn{6}{ c }{Alphabet and Letter Names} \\
        \hline
        a & a & h & hac & q & cu \\
        b & be & i & i & r & ar \\
        c & ce & j & jot & s & es \\
        ç & cedil & k & ka & t & te \\
        d & de & l & el & u & iscon \\
        ð & eð & m & em & v & ve \\
        e & e & n & en & x & ex \\
        f & ef & o & o & y & oy \\
        g & ge & p & pe & z & zet \\
        \hline
    \end{tabular}
\end{table}

The character \orth{w}, which is used in foreign words and some names (like \bl{Willaum} `William'), is pronounced like \orth{v} and called \bl{vescon} `second-vee'. There are also some compound characters. The ligatures \orth{æ œ}, primarily used in Classical loanwords, are called \bl{æsc} and \bl{œthel} respectively and pronounced as if written \orth{e}. The character \orth{ï} is referred to as \bl{i jammel} `twin i' and pronounced as if written \orth{yi}.

In the following we will consider \orth{y} to be a vowel when written between consonants and/or word boundaries, and a consonant otherwise.

\section{Consonants}\label{sec:orth-consonants}
The following letters regularly denote a single phoneme.
\begin{table}[ht]
    \centering
    \begin{tabular}{c c | c c}
        \hline
        b & \phm{b} & m & \phm{m} \\
        ç & \phm{ts} & n & \phm{n} \\
        d & \phm{d} & p & \phm{p} \\
        f & \phm{f} & q & \phm{k} \\
        j & \phm{ʒ} & r & \phm{r} \\
        k & \phm{k} & v & \phm{v} \\
        l & \phm{l} & z & \phm{z} \\
        \hline
    \end{tabular}
\end{table}

Of these, the letters \orth{b d f l m n p r} may appear doubled, which usually does not affect their pronunciation. 
    \glphm{abbað joddy affis ollom commet bannir çoppin arrum}{aˈbaθ ʒɔˈdi aˈfɪz ɔˈlɔm kɔˈmɛt baˈnɪr tsɔˈpɪn aˈrɪm}{abbot cuppa tip relic ban start toffee cosine}
The only exception to this is that \orth{mm nn rr} may denote a single consonant followed by a syllabic consonant. This is the reading when \orth{mm nn rr} occurs word-finally or before a consonant.
    \glphm{yem-m grin-n car-r}{ˈjɛm.m̩ ˈgrɪn.n̩ ˈkar.r̩}{{get-\First\Pl} {judge-\Third\Pl} {seek-\Inf}}
The letter \orth{ð} is pronounced \phm{θ} word-finally and \phm{ð} otherwise.
    \glphm{arð naðr ðeu}{arθ ˈnað.r̩ ðau}{steep viper buttocks}
The letter \orth{x} regularly denotes the phoneme sequence \phm{ks}.
    \glphm{xivol euxon incox}{ksiˈvɔl auˈksɔn ɪnˈkɔks}{waltz quartz digestion}
Consider the final six letters \orth{c g h s t y} when they occur alone (not part of a digraph). In this case, the letters \orth{t y} regularly denote \phm{t j} respectively. The letter \orth{t} may be doubled, without any effect on its pronunciation.
    \glphm{tien yoc | cattin}{tjɛn jɔk | kaˈtɪn}{yours husband | kitten}
The letter \orth{c} is pronounced \phm{ts} before \orth{e i y} and \phm{k} otherwise. The doubled \orth{cc} is pronounced \phm{kts} before \orth{e i y} and \phm{k} otherwise.
    \glphm{cigl cuivr bac | accaç occeir}{tsajl ˈkajv.r̩  bak | aˈkats ɔkˈtsir}{darling copper twig | post kill}
The letter \orth{g} may be pronounced as either of \phm{g ʒ} before \orth{e i y} and is pronounced \phm{g} otherwise. The doubled \orth{gg} may be pronounced as either of \phm{g gʒ} before \orth{e i y} and is pronounced \phm{g} otherwise.
    \glphm{gevou genoil tenguð | suggesson}{geˈvu ʒeˈnɔjl tɛnˈgɪθ | ˌsɪg.ʒɛˈsɔn}{{watch out} knee bauble | suggestion}
The letter \orth{h} is either pronounced \phm{h} or is silent (the latter only in some vocabulary inherited from Latin and after vocalic \orth{g}).
    \glphm{heu hom augher}{hau ɔm oˈjɛr}{mien person tumbler}
The letter \orth{s} is pronounced \phm{z} in coda position or between vowels, and \phm{s} otherwise or when doubled.
    \glphm{bevis bosaç bustr | bossir}{beˈvɪz boˈzats ˈbɪst.r̩  | bɔˈsɪr}{banknote gunfire example | gentleman}

\subsection{Consonant digraphs}

The six letters, \orth{c g h s t y}, combine with other letters to form digraphs. The sequences \orth{sc ch th} regularly denote the phonemes \phm{x k θ}. The potentially ambiguous sequence \orth{sch} is read \phm{sk}.
    \glphm{sclar chym thral | schol}{xlar kɪm θral | skɔl}{lightning lymph slave | school}
The sequence \orth{cq} is pronounced \phm{k}.
    \glphm{jalicq jacquot}{ʒaˈlɪk ʒaˈkwɔt}{waistcoat coat}
The letters \orth{g y} only form digraphs when combined with vowels, for which see below.  

\section{Vowels}\label{sec:orth-vowels}
Divide written vowels into two groups. Free vowels are either word-final or separated from the next vowel in the word by a single consonant letter (or one of \orth{sc ch th}).
    \glphm{choma epu idiocy nothe}{koˈma eˈpi ˌi.djoˈtsi noˈθe}{gnome express stupidity illegitimate}
Checked vowels are either in a final closed syllable or separated from the next vowel in the word by at least two consonant letters.
    \glphm{commun ambasctour sincer}{kɔˈmɪn ˌam.baxˈtur sɪnˈtsɛr}{ordinary emissary genuine}
First consider the letters \orth{a e i o u y} when they occur alone (not part of a digraph). When they are free, \orth{a e i o u y} are pronounced \phm{a e i o i i}.
    \glphm{cla vige bro cru try}{kla viˈʒe bro kri tri}{key wealthy dude raw turn}
When they are checked, \orth{a e i o u y} are pronounced \phm{a ɛ ɪ ɔ ɪ ɪ}.
    \glphm{vindron vampyr effus}{vɪnˈdrɔn vamˈpɪr ɛˈfɪz}{vintner vampire extensive}
Between a consonant and one of the vowels \orth{a e o u}, the letters \orth{i u} denote the semivowel \phm{j}. The letter \orth{u} denotes the semivowel \phm{w} specifically after \orth{q}.
    \glphm{diac suet quelt}{djak sjɛt kwɛlt}{mic familiar dinner}

\subsection{Vowel digraphs}
The digraphs \orth{au eu ou} are pronounced \phm{o au u} respectively. The rarer \orth{ao eo} share their pronunciation with \orth{eu}.
    \glphm{taur teun tout | paon theory}{tor taun tut | paun θauˈri}{bull scant everything | peacock science}
The digraphs \orth{ai ei oi ui} and \orth{ay ey oy uy} are pronounced \phm{e i ɔj aj} respectively. Intervocalically \orth{y} does double duty, simultaneously affecting the previous vowel as stated and acting as an onset \orth{j}.
     \glphm{deit joivr | rayon huyaç}{dit ˈʒɔjv.r̩ | reˈjɔn haˈjats}{finger frost | sunbeam corruption}
The vast majority of coda \orth{g} and many other \orth{g} before consonants form vowel digraphs. Here, the digraphs \orth{ag eg ig og ug aug oug} are usually pronounced \phm{ɛj ɪj aj ɔj aj oj uj} respectively.
     \glphm{bag deg rig pog tug faug voug}{bɛj dɪj raj pɔj taj foj vuj}{berry ten queue few passion sickle wrong}
However, there are exceptions where a \orth{g} in one of the above environments is read consonantally.
    \glphm{magr magr | migrar scigrar}{ˈmag.r̩ ˈmɛj.r̩ | mɪˈgrar xajˈrar}{rash scarce | migrate redistrict}
